\documentclass[a4paper,11pt]{article}

% Encoding and language
\usepackage[utf8]{inputenc}
\usepackage[T1]{fontenc}
\usepackage[brazil]{babel}

% Page layout
\usepackage[margin=2cm, top=1.8cm, bottom=1.8cm]{geometry}

% Math packages
\usepackage{amsmath,amssymb,amsfonts}

% Graphics and TikZ
\usepackage{tikz}
\usetikzlibrary{shapes.geometric, arrows.meta, positioning, fit, backgrounds, calc}

% Colors
\usepackage{xcolor}
\definecolor{checkincolor}{RGB}{66, 133, 244}
\definecolor{poicolor}{RGB}{52, 168, 83}
\definecolor{regioncolor}{RGB}{251, 188, 4}
\definecolor{citycolor}{RGB}{234, 67, 53}
\definecolor{losscolor}{RGB}{128, 0, 128}
\definecolor{codegreen}{RGB}{0, 128, 0}
\definecolor{codered}{RGB}{180, 0, 0}

% Compact sections
\usepackage{titlesec}
\titlespacing*{\section}{0pt}{1.2ex plus 0.3ex minus 0.1ex}{0.6ex plus 0.1ex}
\titlespacing*{\subsection}{0pt}{0.8ex plus 0.2ex minus 0.1ex}{0.4ex plus 0.1ex}
\titleformat{\section}{\normalfont\large\bfseries}{\thesection.}{0.5em}{}
\titleformat{\subsection}{\normalfont\bfseries}{\thesubsection.}{0.5em}{}

% Remove paragraph indent
\setlength{\parindent}{0pt}
\setlength{\parskip}{0.5em}

% Compact lists
\usepackage{enumitem}
\setlist{noitemsep, topsep=2pt, parsep=2pt, partopsep=0pt, leftmargin=1.5em}

% Code listings
\usepackage{listings}
\lstset{
    basicstyle=\ttfamily\small,
    keywordstyle=\color{blue},
    commentstyle=\color{codegreen},
    numbers=none,
    frame=single,
    breaklines=true,
    xleftmargin=0.5em,
    xrightmargin=0.5em
}

\begin{document}

% Title
\begin{center}
{\LARGE\bfseries Check2HGI: Embeddings Hierárquicos de Check-ins via Graph Infomax}\\[0.4em]
{\normalsize Documentação Técnica v2}
\end{center}

\vspace{0.3em}

% Architecture Diagram
\begin{center}
\begin{tikzpicture}[
    node distance=0.8cm and 1.3cm,
    box/.style={rectangle, draw, rounded corners=3pt, minimum width=2.3cm, minimum height=1.1cm, align=center, font=\footnotesize},
    arrow/.style={-{Stealth[length=2mm]}, thick},
    label/.style={font=\scriptsize, text=gray}
]

% Level boxes
\node[box, fill=checkincolor!20, draw=checkincolor] (checkin) {\textbf{Check-ins}\\$[\mathbf{N_c}, D]$};
\node[box, fill=poicolor!20, draw=poicolor, right=of checkin] (poi) {\textbf{POIs}\\$[\mathbf{N_p}, D]$};
\node[box, fill=regioncolor!20, draw=regioncolor, right=of poi] (region) {\textbf{Regiões}\\$[\mathbf{N_r}, D]$};
\node[box, fill=citycolor!20, draw=citycolor, right=of region] (city) {\textbf{Cidade}\\$[\mathbf{D}]$};

% Arrows with labels
\draw[arrow] (checkin) -- node[above, label] {GCN +} node[below, label] {Atenção} (poi);
\draw[arrow] (poi) -- node[above, label] {Atenção +} node[below, label] {GCN} (region);
\draw[arrow] (region) -- node[above, label] {Sigmoid} node[below, label] {Ponderado} (city);

% Loss boundaries
\draw[dashed, losscolor, thick] ($(checkin.south)!0.5!(poi.south) + (0,-0.25)$) -- ++(0,-0.4) node[below, font=\tiny, text=losscolor] {$\mathcal{L}_{c2p}$};
\draw[dashed, losscolor, thick] ($(poi.south)!0.5!(region.south) + (0,-0.25)$) -- ++(0,-0.4) node[below, font=\tiny, text=losscolor] {$\mathcal{L}_{p2r}$};
\draw[dashed, losscolor, thick] ($(region.south)!0.5!(city.south) + (0,-0.25)$) -- ++(0,-0.4) node[below, font=\tiny, text=losscolor] {$\mathcal{L}_{r2c}$};

% Input features
\node[left=0.7cm of checkin, font=\scriptsize, align=right, text width=2cm] (input) {Categoria +\\Temporal};
\draw[arrow, gray] (input) -- (checkin);

\end{tikzpicture}
\end{center}

\vspace{-0.3em}

\section{Visão Geral}

O \textbf{Check2HGI} (Check-in Hierarchical Graph Infomax) é um modelo de embeddings hierárquicos em 4 níveis que aprende representações de check-ins através de agregação progressiva: Check-in $\rightarrow$ POI $\rightarrow$ Região $\rightarrow$ Cidade.

\textbf{Objetivo:} Maximizar a \textbf{informação mútua} entre níveis adjacentes da hierarquia via aprendizado contrastivo auto-supervisionado. O modelo aprende a distinguir pares verdadeiros (check-in pertence ao POI correto) de pares falsos (check-in pareado com POI aleatório), forçando os embeddings a capturar relações semânticas significativas.

\textbf{Arquitetura:}
\begin{itemize}
    \item \textbf{Nível 1 (Check-in):} GCN sobre grafo de sequência temporal do usuário
    \item \textbf{Nível 2 (POI):} Agregação por atenção multi-head dos check-ins
    \item \textbf{Nível 3 (Região):} Agregação por atenção + GCN sobre grafo de adjacência espacial
    \item \textbf{Nível 4 (Cidade):} Agregação ponderada por área geográfica
\end{itemize}

\textbf{Contribuição principal:} Extensão do HGI (3 níveis: POI→Região→Cidade) para incluir o nível de check-in, permitindo embeddings de eventos individuais. Inclui otimização por \textit{corrupção de embeddings} que proporciona speedup de 2x no treinamento.

\textbf{Saídas:} Três conjuntos de embeddings (dimensão $D$, tipicamente 64):
\begin{itemize}
    \item \texttt{embeddings.parquet}: $[N_c, D]$ --- embedding por check-in
    \item \texttt{poi\_embeddings.parquet}: $[N_p, D]$ --- embedding por POI
    \item \texttt{region\_embeddings.parquet}: $[N_r, D]$ --- embedding por região censitária
\end{itemize}

\section{Entrada de Dados}

Cada check-in é representado por um vetor de características que combina:

\textbf{1. Categoria do local} (codificação one-hot): Se existem 50 categorias de POIs, cada check-in tem um vetor de 50 posições com valor 1 na categoria correspondente e 0 nas demais.

\textbf{2. Informação temporal} (codificação cíclica): A hora e o dia da semana são codificados usando seno e cosseno para capturar a natureza cíclica do tempo (23h está próximo de 0h, domingo está próximo de segunda):
\begin{equation*}
\text{temporal} = \left[\sin\left(\frac{2\pi \cdot \text{hora}}{24}\right), \cos\left(\frac{2\pi \cdot \text{hora}}{24}\right), \sin\left(\frac{2\pi \cdot \text{dia}}{7}\right), \cos\left(\frac{2\pi \cdot \text{dia}}{7}\right)\right]
\end{equation*}

\textbf{Grafo de check-ins:} Os check-ins são conectados em um grafo onde arestas ligam visitas consecutivas do mesmo usuário. O peso de cada aresta diminui exponencialmente com o tempo entre visitas:
\begin{equation*}
w_{ij} = \exp\left(-\frac{\text{tempo entre } i \text{ e } j}{\tau}\right), \quad \tau = 3600\text{s (1 hora)}
\end{equation*}

Isso significa que check-ins próximos no tempo têm conexões mais fortes.

\section{Codificadores (Encoders)}

\textbf{Nível 1 - CheckinEncoder:} Usa uma Rede Convolucional de Grafos (GCN) que propaga informações entre check-ins conectados. Cada check-in ``aprende'' com seus vizinhos no grafo:
\begin{equation*}
\mathbf{H}^{(l+1)} = \text{PReLU}\left(\tilde{\mathbf{D}}^{-1/2}\tilde{\mathbf{A}}\tilde{\mathbf{D}}^{-1/2}\mathbf{H}^{(l)}\mathbf{W}^{(l)}\right)
\end{equation*}

\textbf{Nível 2 - Checkin2POI:} Agrega todos os check-ins de um POI usando atenção multi-head. O mecanismo de atenção aprende quais check-ins são mais importantes para caracterizar cada POI:
\begin{equation*}
\mathbf{z}_p = \sum_{i \in \mathcal{C}_p} \alpha_i \mathbf{v}_i, \quad \text{onde } \alpha_i = \text{softmax}\left(\frac{\mathbf{q}^\top \mathbf{k}_i}{\sqrt{d}}\right)
\end{equation*}

\textbf{Nível 3 - POI2Region:} Similar ao nível anterior, agrega POIs em regiões e aplica GCN no grafo de adjacência regional (regiões vizinhas geograficamente).

\textbf{Nível 4 - Region2City:} Média ponderada das regiões pela sua área geográfica, passada por função sigmoid.

%--- PAGE BREAK ---%
\newpage

\section{Otimização por Corrupção de Embeddings}

Esta seção explica uma otimização importante que dobra a velocidade de treinamento.

\subsection{Aprendizado Contrastivo: A Ideia Básica}

O modelo aprende comparando pares ``positivos'' (corretos) com pares ``negativos'' (incorretos). Por exemplo, para a fronteira Check-in↔POI:
\begin{itemize}
    \item \textbf{Par positivo:} embedding do check-in $c_1$ + embedding do POI onde $c_1$ ocorreu
    \item \textbf{Par negativo:} embedding do check-in $c_1$ + embedding de um POI aleatório
\end{itemize}

O modelo deve aprender a dar pontuação alta para pares positivos e baixa para negativos.

\subsection{Método Original: Corrupção de Features}

No método tradicional (usado no Deep Graph Infomax original), os pares negativos são criados ``corrompendo'' as características de entrada:

\begin{lstlisting}[language=Python, title={\textbf{Método Original (2 passagens pelo encoder)}}]
# Passagem 1: dados reais
pos_embeddings = encoder(features, graph)

# Passagem 2: dados corrompidos (features embaralhadas)
corrupted_features = shuffle(features)  # mistura features entre nos
neg_embeddings = encoder(corrupted_features, graph)

# Calcula perda comparando pos vs neg
loss = contrastive_loss(pos_embeddings, neg_embeddings)
\end{lstlisting}

\textbf{Problema:} O encoder (a rede GCN) precisa processar os dados \textbf{duas vezes} --- uma para os dados reais e outra para os corrompidos. Como o encoder é a parte mais custosa computacionalmente, isso dobra o tempo de treinamento.

\subsection{Método Otimizado: Corrupção de Embeddings}

A insight chave é: para o aprendizado contrastivo funcionar, os pares negativos só precisam ser ``errados'' --- não importa \textit{como} foram gerados. Então, ao invés de corromper a entrada e rodar o encoder de novo, podemos simplesmente embaralhar os embeddings de saída:

\begin{lstlisting}[language=Python, title={\textbf{Método Otimizado (1 passagem pelo encoder)}}]
# Passagem unica: dados reais
pos_embeddings = encoder(features, graph)

# Negativos: embaralha os embeddings ja calculados
neg_embeddings = pos_embeddings[random_permutation]

# Calcula perda comparando pos vs neg
loss = contrastive_loss(pos_embeddings, neg_embeddings)
\end{lstlisting}

\textbf{Por que funciona?} O objetivo do aprendizado contrastivo é ensinar o modelo que ``check-in $c_1$ pertence ao POI $p_1$''. Um par negativo válido é qualquer combinação incorreta. Ao embaralhar os embeddings, garantimos que cada check-in será pareado com o POI errado --- exatamente o que precisamos.

\textbf{Resultado:} \colorbox{yellow!30}{\textbf{Speedup de 2x}} no treinamento, pois o encoder (operação mais cara) roda apenas uma vez por época.

\section{Função de Perda}

A perda total combina três termos, um para cada fronteira hierárquica:
\begin{equation*}
\mathcal{L} = \underbrace{0.4 \cdot \mathcal{L}_{c2p}}_{\text{Check-in}\leftrightarrow\text{POI}} + \underbrace{0.3 \cdot \mathcal{L}_{p2r}}_{\text{POI}\leftrightarrow\text{Região}} + \underbrace{0.3 \cdot \mathcal{L}_{r2c}}_{\text{Região}\leftrightarrow\text{Cidade}}
\end{equation*}

Cada termo usa um \textbf{discriminador bilinear} que calcula a ``compatibilidade'' entre dois embeddings:
\begin{equation*}
\mathcal{D}(\mathbf{e}_1, \mathbf{e}_2) = \sigma(\mathbf{e}_1^\top \mathbf{W} \mathbf{e}_2)
\end{equation*}

onde $\mathbf{W}$ é uma matriz de pesos aprendida e $\sigma$ é a função sigmoid (saída entre 0 e 1).

A perda para cada fronteira maximiza a pontuação de pares corretos e minimiza a de incorretos:
\begin{equation*}
\mathcal{L}_{*} = -\underbrace{\log \mathcal{D}(\mathbf{e}^+, \mathbf{e}^+)}_{\text{maximizar para pares corretos}} - \underbrace{\log(1 - \mathcal{D}(\mathbf{e}^+, \mathbf{e}^-))}_{\text{minimizar para pares incorretos}}
\end{equation*}

\textbf{Intuição:} É como treinar um classificador binário que deve responder ``sim'' para pares verdadeiros e ``não'' para falsos. Se o modelo consegue fazer isso bem, significa que os embeddings capturam informações significativas sobre a hierarquia espacial.

\section{Referências}

\begin{enumerate}[label={[\arabic*]}, leftmargin=2em, itemsep=0.2em]
\item Veličković, P. et al. ``Deep Graph Infomax''. \textit{ICLR}, 2019. --- Base teórica para maximização de informação mútua em grafos.
\item Lee, J. et al. ``Set Transformer: A Framework for Attention-based Permutation-Invariant Neural Networks''. \textit{ICML}, 2019. --- Mecanismo de atenção para agregação.
\item Kipf, T. \& Welling, M. ``Semi-Supervised Classification with Graph Convolutional Networks''. \textit{ICLR}, 2017. --- Redes convolucionais de grafos.
\end{enumerate}

\end{document}
